\documentclass[border=4pt]{standalone}
\usepackage{amsmath,amssymb,mathtools,xcolor,varwidth}
\begin{document}
\begin{varwidth}{28em}
\large\bfseries
Now shrink everything: let $B,D,G$ approach $A$. The intersection $J$ of $BG$ and      $AG$ races toward $A$ so that $GJ$ becomes less than any assignable      length, whence $AG:Ag$ differs from equality by nothing in the limit. With      $AB^{2}=AG\times BD$ and $Ab^{2}=Ag\times bd$, the ratio $AB^{2}:Ab^{2}$      is compounded of $AG:Ag$ and $BD:bd$; since the first factor becomes unity, by      Lemma I the ultimate ratio $AB^{2}:Ab^{2}$ equals $BD:bd$. Thus the subtense      goes as the \textit{square} of the arc — every smooth bend, circle or      parabola alike, peels from its tangent quadratically, the universal      signature of finite curvature. \textit{Q.E.D.}
\end{varwidth}
\end{document}
